\begin{tikzpicture}[
  font=\scriptsize,
  execute at begin node={\hyphenpenalty=10000\exhyphenpenalty=10000},
  kotak/.style={draw, rounded corners, fill=blue!8, text width=2.45cm, minimum height=1.45cm, align=center},
  tekan/.style={kotak, fill=orange!12},
  hasil/.style={kotak, fill=green!10},
  masalah/.style={kotak, fill=red!22},
  label/.style={font=\tiny, inner sep=1pt},
  panah/.style={-{Stealth[length=2mm]}, thick}
]
  % Tiga gerbang diperiksa berurutan. Gerbang ketiga butuh bukti dari konsumen.
  \node[kotak] (usul) at (0,0) {\textbf{Usulan perubahan}\\{\tiny\mdseries satu butir kontrak}};
  \node[tekan] (g1) at (3.1,0) {\textbf{Gerbang 1}\\Janji tertulis dihapus, dipersempit, atau artinya berubah?};
  \node[tekan] (g2) at (6.2,0) {\textbf{Gerbang 2}\\Ada hal baru pada jawaban, atau perilaku teramati berubah?};
  \node[tekan] (g3) at (9.3,0) {\textbf{Gerbang 3}\\Uji konsumen: ada konsumen yang bergantung?};
  \node[masalah] (rusak) at (3.1,-2.75) {\textbf{Merusak}\\Versi utama baru, jalan berdampingan, masa transisi di Bab 10};
  \node[hasil] (aman) at (6.2,-2.75) {\textbf{Aman}\\Versi minor, konsumen tetap diberi tahu};
  \node[masalah] (rusak2) at (9.3,-2.75) {\textbf{Merusak bagi konsumen itu}\\Ubah usulan atau jalankan migrasi};

  \draw[panah] (usul) -- (g1);
  \draw[panah] (g1) -- node[label, above] {tidak} (g2);
  \draw[panah] (g2) -- node[label, above] {ya} (g3);
  \draw[panah] (g1) -- node[label, right] {ya} (rusak);
  \draw[panah] (g2) -- node[label, right] {tidak} (aman);
  \draw[panah] ([xshift=0.45cm]g3.south) -- node[label, right] {ya} ([xshift=0.45cm]rusak2.north);
  \draw[panah] ([xshift=-0.75cm]g3.south) -- node[label, left, pos=0.4] {tidak} (aman.north east);
\end{tikzpicture}
